\subsection{Entorno experimental}
\label{sec:experimental-environment}

Todos los experimentos fueron ejecutados bajo el mismo en   torno de hardware y software.
La Tabla~\ref{tab:environment} documenta la configuración utilizada con el objetivo de
facilitar la reproducibilidad de los resultados.

\begin{table}[H]
\centering
\caption{Entorno utilizado para ejecutar los experimentos.}
\label{tab:environment}
\begin{tabularx}{0.9\textwidth}{lY}
\toprule
\textbf{Elemento} & \textbf{Configuración} \\
\midrule

Procesador &
\todoresult{modelo exacto de CPU} \\

Arquitectura &
\todoresult{x86\_64, ARM64, etc.} \\

Núcleos físicos &
\todoresult{valor} \\

Núcleos lógicos &
\todoresult{valor} \\

Memoria RAM &
\todoresult{capacidad total} \\

Sistema operativo &
\todoresult{distribución y versión} \\

Kernel &
\todoresult{versión} \\

Docker &
\todoresult{versión} \\

Imagen base &
\todoresult{imagen y versión} \\

SageMath &
\todoresult{versión} \\

Python &
\todoresult{versión} \\

Biblioteca criptográfica &
\todoresult{nombre y versión} \\

Código del experimento &
\todoresult{commit Git evaluado} \\

Fecha de ejecución &
\todoresult{fecha} \\

\bottomrule
\end{tabularx}
\end{table}

Todos los candidatos fueron ejecutados sobre el mismo hardware y utilizando el mismo
mecanismo de instrumentación. Durante la toma principal de resultados se evitó ejecutar
aplicaciones ajenas al experimento que pudieran introducir carga significativa sobre
el sistema.


% ============================================================
\subsection{Configuración criptográfica}
\label{sec:crypto-configuration}

Los parámetros criptográficos empleados para cada candidato se presentan en la
Tabla~\ref{tab:crypto-parameters}.

\begin{table}[H]
\centering
\caption{Parámetros criptográficos utilizados en cada implementación.}
\label{tab:crypto-parameters}
\small
\begin{tabularx}{\textwidth}{lYYYY}
\toprule
\textbf{Protocolo} &
\textbf{Dimensión} &
\textbf{Módulo / curva} &
\textbf{Distribución / error} &
\textbf{Otros parámetros} \\
\midrule

ECDSA &
N/A &
\todoresult{curva} &
N/A &
\todoresult{hash, codificación, etc.} \\

LWE estándar &
\todoresult{$n$, $m$} &
\todoresult{$q$} &
\todoresult{$\chi_s$, $\chi_e$, $\sigma$} &
\todoresult{otros} \\

Binary-LWE &
\todoresult{$n$, $m$} &
\todoresult{$q$} &
\todoresult{secreto y error} &
\todoresult{otros} \\

Ring-LWE &
\todoresult{$n$} &
\todoresult{$q$, anillo} &
\todoresult{distribuciones} &
\todoresult{polinomio/modulus} \\

LWR &
\todoresult{$n$, $m$} &
\todoresult{$q$, $p$} &
\todoresult{distribuciones} &
\todoresult{redondeo} \\

LWE propuesto &
\todoresult{parámetros} &
\todoresult{parámetros} &
\todoresult{parámetros} &
\todoresult{parámetros} \\

\bottomrule
\end{tabularx}
\end{table}

Los parámetros fueron seleccionados buscando un nivel de seguridad comparable entre
los candidatos. Cuando no fue posible establecer una equivalencia exacta, esta
limitación se indica explícitamente en la Sección~\ref{sec:security-evaluation}.


% ============================================================
\subsection{Diseño de los experimentos}
\label{sec:experimental-design}

\subsubsection{Número de ejecuciones}

Cada operación fue ejecutada:

\begin{equation}
    N = \todoresult{número de iteraciones}
\end{equation}

veces por configuración.

Antes de iniciar la toma de datos se realizaron:

\begin{equation}
    N_{\text{warm-up}} =
    \todoresult{número de iteraciones de calentamiento}
\end{equation}

ejecuciones de calentamiento que no fueron incluidas en los resultados finales.

\subsubsection{Estadísticos reportados}

Para evitar representar el rendimiento mediante una única medición, para cada operación
se calcularon los siguientes estadísticos:

\begin{itemize}
    \item media aritmética;
    \item mediana;
    \item desviación estándar;
    \item mínimo;
    \item máximo;
    \item percentil 95;
    \item percentil 99.
\end{itemize}



% ============================================================
\subsection{Métricas evaluadas}
\label{sec:metrics}

Las métricas utilizadas se agrupan en siete dimensiones principales:

\begin{enumerate}
    \item rendimiento computacional;
    \item consumo de recursos;
    \item comunicación;
    \item seguridad;
    \item complejidad de implementación;
    \item características operacionales y experiencia de usuario;
    \item madurez y facilidad de despliegue.
\end{enumerate}


% ============================================================
\subsubsection{Rendimiento computacional}

Para cada protocolo se mide el tiempo requerido por:

\begin{itemize}
    \item generación de claves;
    \item generación del reto;
    \item generación de la respuesta;
    \item verificación;
    \item autenticación completa.
\end{itemize}

Las mediciones se reportan en milisegundos.

Adicionalmente se calcula el throughput:

\begin{equation}
    R =
    \frac{N_{\text{auth}}}
         {T_{\text{total}}},
\end{equation}

expresado en autenticaciones por segundo.


% ============================================================
\subsubsection{Consumo de memoria y almacenamiento}

Para cada implementación se estudian:

\begin{itemize}
    \item memoria RAM promedio;
    \item memoria RAM máxima;
    \item tamaño de la clave pública;
    \item tamaño de la clave privada;
    \item tamaño de parámetros públicos adicionales;
    \item estado temporal requerido durante una autenticación;
    \item almacenamiento persistente requerido por usuario.
\end{itemize}

Los tamaños de estructuras criptográficas se calculan utilizando su representación
serializada real y no únicamente su representación matemática abstracta.


% ============================================================
\subsubsection{Comunicación}

El volumen de comunicación se calcula utilizando los bytes efectivamente serializados
por la aplicación.

Para una autenticación:

\begin{equation}
    B_{\text{auth}}
    =
    B_{\text{request}}
    +
    B_{\text{challenge}}
    +
    B_{\text{response}}
    +
    B_{\text{metadata}}.
\end{equation}

Además se registra:

\begin{itemize}
    \item número de mensajes intercambiados;
    \item número de viajes de ida y vuelta (\textit{round trips});
    \item tamaño máximo de un mensaje;
    \item cantidad total de bytes transmitidos.
\end{itemize}


% ============================================================
\subsubsection{Seguridad}

La evaluación de seguridad considera:

\begin{itemize}
    \item problema matemático subyacente;
    \item parámetros utilizados;
    \item estimación de seguridad clásica;
    \item estimación de seguridad frente a adversarios cuánticos;
    \item ataques conocidos relevantes;
    \item sensibilidad a una selección incorrecta de parámetros;
    \item requisitos de aleatoriedad;
    \item probabilidad de fallo cuando sea aplicable;
    \item vulnerabilidades derivadas de la implementación.
\end{itemize}


% ============================================================
\subsubsection{Complejidad de implementación}

La dificultad de implementación no se asigna únicamente mediante una valoración
subjetiva. Se consideran los siguientes factores:

\begin{itemize}
    \item cantidad de primitivas criptográficas necesarias;
    \item cantidad de parámetros configurables;
    \item líneas de código asociadas con la implementación criptográfica;
    \item número y naturaleza de dependencias externas;
    \item necesidad de aritmética especializada;
    \item complejidad del muestreo aleatorio;
    \item complejidad de serialización;
    \item manejo de estado;
    \item disponibilidad de bibliotecas consolidadas.
\end{itemize}


% ============================================================
\subsubsection{Características operacionales}

Para estudiar las implicaciones en un sistema real se consideran:

\begin{itemize}
    \item número de pasos requeridos durante el registro;
    \item número de pasos requeridos durante una autenticación;
    \item necesidad de almacenar secretos en el dispositivo del usuario;
    \item necesidad de servidor en línea;
    \item posibilidad de verificación fuera de línea;
    \item necesidad de mantener estado por usuario;
    \item procedimiento de renovación de claves;
    \item procedimiento ante compromiso de credenciales;
    \item requerimientos de infraestructura adicionales.
\end{itemize}
